\documentclass[11pt]{article}
% Shared preamble for manuscripts-v2 papers. Load with % Shared preamble for manuscripts-v2 papers. Load with % Shared preamble for manuscripts-v2 papers. Load with \input{../common/preamble.tex}
% immediately after \documentclass[11pt]{article}.
\ifdefined\pdftrailerid\pdftrailerid{}\fi
\usepackage[T1]{fontenc}
\usepackage{lmodern}
\usepackage[letterpaper,margin=1.02in]{geometry}
\usepackage{amsmath,amssymb,amsthm,mathtools,booktabs,array,microtype,enumitem}
\usepackage{tikz}
\usepackage{xurl}
\usepackage[colorlinks=true,allcolors=bluegreen]{hyperref}
\definecolor{bluegreen}{RGB}{22,83,91}
\usepackage[capitalize,nameinlink]{cleveref}

\theoremstyle{plain}
\newtheorem{theorem}{Theorem}[section]
\newtheorem{lemma}[theorem]{Lemma}
\newtheorem{proposition}[theorem]{Proposition}
\newtheorem{corollary}[theorem]{Corollary}
\newtheorem{conjecture}[theorem]{Conjecture}
\newtheorem*{theorem*}{Theorem}
\newtheorem{maintheorem}{Theorem}
\renewcommand{\themaintheorem}{\Alph{maintheorem}}
\theoremstyle{definition}
\newtheorem{definition}[theorem]{Definition}
\newtheorem{example}[theorem]{Example}
\newtheorem{question}[theorem]{Question}
\theoremstyle{remark}
\newtheorem{remark}[theorem]{Remark}
\crefname{maintheorem}{Theorem}{Theorems}
\crefname{conjecture}{Conjecture}{Conjectures}
\crefname{question}{Question}{Questions}
\crefname{example}{Example}{Examples}

\newcommand{\R}{\mathbb R}
\newcommand{\Q}{\mathbb Q}
\newcommand{\Z}{\mathbb Z}
\newcommand{\N}{\mathbb N}
\newcommand{\T}{\mathbb T}
\newcommand{\C}{\mathbb C}
\DeclareMathOperator{\den}{den}
\newcommand{\e}{\mathrm e}

% Immutable Lean source links at research-checkpoint-18.
\newcommand{\leancommit}{e4809cf561573aef588e6ffc491a35cdb79fd9ec}
\newcommand{\leanbase}{https://github.com/eldernos/MathematicsResearch/blob/\leancommit}
% \lean{File}{line}{label}: link to NormalityChecker/File.lean#Lline at the pinned commit.
\newcommand{\lean}[3]{\href{\leanbase/NormalityChecker/#1.lean\#L#2}{#3}}
% \repo{path}{label}: link to any repository path at the pinned commit.
\newcommand{\repo}[2]{\href{https://github.com/eldernos/MathematicsResearch/tree/\leancommit/#1}{#2}}

% Evidence labels (see papers/WAVE2-CONTRACT.md).
\newcommand{\evidence}[1]{\textup{[\emph{#1}]}}

\setlength{\emergencystretch}{1.5em}
\raggedbottom

% immediately after \documentclass[11pt]{article}.
\ifdefined\pdftrailerid\pdftrailerid{}\fi
\usepackage[T1]{fontenc}
\usepackage{lmodern}
\usepackage[letterpaper,margin=1.02in]{geometry}
\usepackage{amsmath,amssymb,amsthm,mathtools,booktabs,array,microtype,enumitem}
\usepackage{tikz}
\usepackage{xurl}
\usepackage[colorlinks=true,allcolors=bluegreen]{hyperref}
\definecolor{bluegreen}{RGB}{22,83,91}
\usepackage[capitalize,nameinlink]{cleveref}

\theoremstyle{plain}
\newtheorem{theorem}{Theorem}[section]
\newtheorem{lemma}[theorem]{Lemma}
\newtheorem{proposition}[theorem]{Proposition}
\newtheorem{corollary}[theorem]{Corollary}
\newtheorem{conjecture}[theorem]{Conjecture}
\newtheorem*{theorem*}{Theorem}
\newtheorem{maintheorem}{Theorem}
\renewcommand{\themaintheorem}{\Alph{maintheorem}}
\theoremstyle{definition}
\newtheorem{definition}[theorem]{Definition}
\newtheorem{example}[theorem]{Example}
\newtheorem{question}[theorem]{Question}
\theoremstyle{remark}
\newtheorem{remark}[theorem]{Remark}
\crefname{maintheorem}{Theorem}{Theorems}
\crefname{conjecture}{Conjecture}{Conjectures}
\crefname{question}{Question}{Questions}
\crefname{example}{Example}{Examples}

\newcommand{\R}{\mathbb R}
\newcommand{\Q}{\mathbb Q}
\newcommand{\Z}{\mathbb Z}
\newcommand{\N}{\mathbb N}
\newcommand{\T}{\mathbb T}
\newcommand{\C}{\mathbb C}
\DeclareMathOperator{\den}{den}
\newcommand{\e}{\mathrm e}

% Immutable Lean source links at research-checkpoint-18.
\newcommand{\leancommit}{e4809cf561573aef588e6ffc491a35cdb79fd9ec}
\newcommand{\leanbase}{https://github.com/eldernos/MathematicsResearch/blob/\leancommit}
% \lean{File}{line}{label}: link to NormalityChecker/File.lean#Lline at the pinned commit.
\newcommand{\lean}[3]{\href{\leanbase/NormalityChecker/#1.lean\#L#2}{#3}}
% \repo{path}{label}: link to any repository path at the pinned commit.
\newcommand{\repo}[2]{\href{https://github.com/eldernos/MathematicsResearch/tree/\leancommit/#1}{#2}}

% Evidence labels (see papers/WAVE2-CONTRACT.md).
\newcommand{\evidence}[1]{\textup{[\emph{#1}]}}

\setlength{\emergencystretch}{1.5em}
\raggedbottom

% immediately after \documentclass[11pt]{article}.
\ifdefined\pdftrailerid\pdftrailerid{}\fi
\usepackage[T1]{fontenc}
\usepackage{lmodern}
\usepackage[letterpaper,margin=1.02in]{geometry}
\usepackage{amsmath,amssymb,amsthm,mathtools,booktabs,array,microtype,enumitem}
\usepackage{tikz}
\usepackage{xurl}
\usepackage[colorlinks=true,allcolors=bluegreen]{hyperref}
\definecolor{bluegreen}{RGB}{22,83,91}
\usepackage[capitalize,nameinlink]{cleveref}

\theoremstyle{plain}
\newtheorem{theorem}{Theorem}[section]
\newtheorem{lemma}[theorem]{Lemma}
\newtheorem{proposition}[theorem]{Proposition}
\newtheorem{corollary}[theorem]{Corollary}
\newtheorem{conjecture}[theorem]{Conjecture}
\newtheorem*{theorem*}{Theorem}
\newtheorem{maintheorem}{Theorem}
\renewcommand{\themaintheorem}{\Alph{maintheorem}}
\theoremstyle{definition}
\newtheorem{definition}[theorem]{Definition}
\newtheorem{example}[theorem]{Example}
\newtheorem{question}[theorem]{Question}
\theoremstyle{remark}
\newtheorem{remark}[theorem]{Remark}
\crefname{maintheorem}{Theorem}{Theorems}
\crefname{conjecture}{Conjecture}{Conjectures}
\crefname{question}{Question}{Questions}
\crefname{example}{Example}{Examples}

\newcommand{\R}{\mathbb R}
\newcommand{\Q}{\mathbb Q}
\newcommand{\Z}{\mathbb Z}
\newcommand{\N}{\mathbb N}
\newcommand{\T}{\mathbb T}
\newcommand{\C}{\mathbb C}
\DeclareMathOperator{\den}{den}
\newcommand{\e}{\mathrm e}

% Immutable Lean source links at research-checkpoint-18.
\newcommand{\leancommit}{e4809cf561573aef588e6ffc491a35cdb79fd9ec}
\newcommand{\leanbase}{https://github.com/eldernos/MathematicsResearch/blob/\leancommit}
% \lean{File}{line}{label}: link to NormalityChecker/File.lean#Lline at the pinned commit.
\newcommand{\lean}[3]{\href{\leanbase/NormalityChecker/#1.lean\#L#2}{#3}}
% \repo{path}{label}: link to any repository path at the pinned commit.
\newcommand{\repo}[2]{\href{https://github.com/eldernos/MathematicsResearch/tree/\leancommit/#1}{#2}}

% Evidence labels (see papers/WAVE2-CONTRACT.md).
\newcommand{\evidence}[1]{\textup{[\emph{#1}]}}

\setlength{\emergencystretch}{1.5em}
\raggedbottom

% Author information approved by the repository owner on 2026-11-25:
%   name: Christopher Anderson; designation: Independent Researcher.
% No academic affiliation, ORCID or personal contact details have been supplied.
% Do not add any that the owner has not approved.
\newcommand{\OwnerAuthor}{Christopher Anderson}
\newcommand{\OwnerAffiliation}{Independent Researcher}
\newcommand{\OwnerContact}{Corrections and comments: \url{https://github.com/eldernos/MathematicsResearch/issues}}
\newcommand{\ReleaseLabel}{manuscripts-v2.0.0}
\newcommand{\LicenseLine}{\copyright\ 2026 Christopher Anderson. Licensed under \href{https://creativecommons.org/licenses/by/4.0/}{CC BY 4.0}.}
% \makebyline{version}{date}: centred author block used under the title.
\newcommand{\makebyline}[2]{%
  {\large \OwnerAuthor\par}%
  \ifx\OwnerAffiliation\empty\else{\small \OwnerAffiliation\par}\fi
  \vspace{0.3em}
  {\small Version #1 \enspace$\cdot$\enspace #2 \enspace$\cdot$\enspace \ReleaseLabel\par}%
  {\footnotesize \OwnerContact\par}%
  {\footnotesize \LicenseLine\par}%
  \hypersetup{pdfauthor={\OwnerAuthor}}}

\hypersetup{pdftitle={Polynomial Images of Stoneham Numbers}}
\newcommand{\ee}[1]{\exp(2\pi i #1)}
\begin{document}
\begin{center}
{\LARGE Polynomial Images of Stoneham Numbers\par}
\medskip
\makebyline{2.0}{25 November 2026}
\end{center}
\begin{abstract}
Mira Cult's S08 (Theorem 1.1 and Corollaries 1.2--1.3) proves polynomial
normality and joint distribution for Stoneham series with a common good prime.
We give a complete kernel-verified proof for the ordinary series with coprime
radix \(b\ge2\) and parameter \(c>2\). Inclusion--exclusion bounds cancellation
in its leading denominator. S08, Theorem 1.1(b), already gives an effective power saving;
we specify a scalar discrepancy exponent in each degree.
For \(c=2\) we determine the square's prefix denominators and give a
normality criterion for quadratics. A written transfer from exact finite
certificates proves normality of every rational polynomial of degree
\(3\) through \(120\), of degree \(128\), or of degree \(ep^r\) for sixteen
pairs \((e,p)\) and \(r\ge0\). We determine the least valuation of a
permutation-orbit contribution in every degree. Normality of the
resonant quadratic \(\alpha_{b,2}^2-2\alpha_{b,2}\) remains open.
\end{abstract}

\section{Polynomial values and simultaneous digits}
Theorems A and B are the ordinary case of Mira Cult's public, unrefereed
S08 \cite[Thm.~1.1 and Cors.~1.2--1.3]{mira2}. We give a kernel-verified
proof using inclusion--exclusion on upper sets, then treat the boundary
case \(c=2\) and compute scalar discrepancy exponents.
For coprime integers \(b,c\ge2\), let
\[
 \alpha_{b,c}=\sum_{j=1}^{\infty}\frac{1}{c^j b^{c^j}}.
\]
A real number \(x\) is normal in base \(b\) if every word of length \(m\ge1\)
has limiting frequency \(b^{-m}\) in its fractional digits. Equivalently,
the fractional parts \(\{b^nx\}\), \(n=0,1,\ldots\), are uniformly distributed.
\begin{maintheorem}[Mira Cult, S08, Theorem 1.1(a) and Corollary 1.2, \(q=c\)]\label{main}
If \(c>2\), then \(P(\alpha_{b,c})\) is normal in base \(b\) for every
nonconstant \(P\in\Q[X]\). \evidence{Kernel-verified}
\end{maintheorem}
\begin{maintheorem}[Mira Cult, S08, Corollaries 1.2 and 1.3]\label{joint}
Under the same hypotheses let \(r_1,\ldots,r_m\) be distinct positive integers.
For \(0\le u_i\le v_i\le1\) and \(\alpha=\alpha_{b,c}\),
\[
 \lim_{N\to\infty}\frac1N\#\{0\le n<N:\ \{b^n\alpha^{r_i}\}\in[u_i,v_i)
             \text{ for all }i\}=\prod_{i=1}^m(v_i-u_i).
\]
\evidence{Kernel-verified}
\end{maintheorem}
The \lean{PeriodicPolynomialNormal}{145}{polynomial theorem} and
\lean{PeriodicPolynomialNormal}{170}{rectangle theorem} use the same
starting digit in every coordinate. Transcendence of the generator follows
already from Coons's Mahler-method argument \cite[Cor.~A.3 of the 2013 preprint]{coons}.
It also follows from Theorem~\ref{main}, since a polynomial relation would
produce the nonnormal value zero.

\paragraph{Antecedents.}
Stoneham's binary example \cite{stoneham}, A.~N.~Korobov's lacunary-series
criterion \cite{korobov,kano} and Bailey--Crandall's coprime-parameter
theorem \cite[Cor.~4.9]{bc} account for the generator itself.
Bugeaud surveys the discrepancy of its radix orbit \cite[\S5.1]{bugeaud}.
Kano--Shiokawa and Wagner constructed rings with normal polynomial values
\cite{ks,wagner}. The hypotheses of \cite[Thms.~4--5, institutional preprint (15)--(16)]{ks} require two
growth ratios to diverge. For the standard exponents
\(\lambda_j=j,\ \mu_j=c^j\), these ratios are \(j/c^{j-1}\) and \(\log c\).
Nakai--Shiokawa's construction concatenates polynomial values at primes
\cite{ns}. Here the polynomial is evaluated on a single real number.
N.~M.~Korobov studied synchronized digit distributions \cite{jointkorobov}.

Mira Cult's archive of 7 October 2026 \cite[S08, Thm.~1.1 and Cors.~1.2--1.3]{mira2} states polynomial
normality for \(\sum_{j\ge1}q^{-j}b^{-c^j}\) when \(q>2\) and some prime
divides both \(c\) and \(q\) but does not divide \(b\). It also gives a
power saving for the synchronized vector of the first \(m\) powers.
Setting \(q=c\) gives Theorem~\ref{main}, and coordinate projection gives
Theorem~\ref{joint}. More generally, S08, Corollary 1.3 gives joint
normality precisely when the polynomial classes in \(\Q[X]/\Q\) are
linearly independent: every nonzero integer projection then has positive
degree, whereas a rational constant projection has finite radix orbit. S08 obtains its monic polynomial by separating high
and low indices outside short intervals following insertions. Here
inclusion--exclusion gives a monic polynomial on every reflected upper set.
Both arguments evaluate at a nonintegral algebraic tail. The notes
of 23 September 2026 \cite[S07, S09]{mira} contain the earlier
single-series and joint prime/order results. The explicit scalar rates below
are deductions from Vandehey's bounds \cite{vandehey}. Normality here
concerns the defining radix. Bailey--Borwein \cite[Thm.~2]{bb}
establish nonnormality of generators in certain other radices. For the
square of \(\alpha_{2,3}\), the same value is normal in base two and
\lean{AlternateBaseStoneham}{385}{nonnormal in base 48}
\evidence{Kernel-verified}. Degree one in the
discrepancy estimate is also classical: Bailey--Crandall give the rate
after their Theorem~4.8, and Bugeaud records it for the real orbit
\cite[p.~535]{bc}, \cite[\S5.1]{bugeaud}.

\subsection*{Notation for digit positions and rational prefixes}
The \(k\)th \emph{shell} consists of positions \(c^k\le n<c^{k+1}\).
A property holds on late shells if it holds for all sufficiently large
\(k\). For a tuple \(j=(j_1,\ldots,j_d)\) of positive indices write
\(\sigma(j)=\sum_i c^{j_i}\). Its \emph{insertion} occurs at the first
digit position \(n=\sigma(j)\) at which all its radix factors have been
shifted out of the denominator. The shifted prefix \(R_d(n)\) is the
rational number obtained from \(b^n\alpha^d\) by retaining exactly the
tuples with \(\sigma(j)\le n\):
\[
 R_d(n)=\sum_{\substack{j_i\ge1\\\sigma(j)\le n}}
             c^{-\sum_i j_i} b^{n-\sigma(j)},\qquad R_0(n)=b^n .
\]
Absolute convergence justifies the expansion. For \(P=\sum_\ell P_\ell X^\ell\)
put \(R_P(n)=\sum_\ell P_\ell R_\ell(n)\). The \emph{deficit}
\(v_i=k-j_i\) counts how far an index lies below the largest available
index \(k\). With \(t=n/c^k\), reflection gives the \emph{weak cut}
\[
 U_t=\{v\in\N^d:\ \textstyle\sum_i c^{-v_i}\le t\}.
\]
Here \(0\in\N\); the prefix uses only \(v_i<k\). The \emph{strict cut}
uses \(<t\). Both are upper sets: increasing coordinates preserves
membership. Every weak shell cut contains the fixed vertex
\((J,\ldots,J)\) for \(c^J\ge d\); we call this vertex the anchor.
Write \(|v|=\sum_i v_i\). For a prime \(p\mid c\) and periodic unit
weights \(u_v\) in \(\Z_p\), the \emph{mass} of an upper set is
\[
 H(U)=\sum_{v\in U}\prod_i c^{v_i}u_{v_i}\quad\text{in }\Z_p.
\]
Only finitely many summands survive modulo any \(p^L\), so the series
converges. Its \emph{loss} is \(v_p(H(U))\), with \(v_p(0)=+\infty\):
each unit of loss cancels one power of \(p\) from the nominal denominator.

\section{A nonintegral sum controls the denominator}\label{mass}
\begin{proposition}\label{massprop}
For \(c>2\), a periodic sequence of roots of unity in a number field
embedded in \(\Q_p\), and \(p\mid c\), the convergent \(p\)-adic mass \(H(U)\) of every nonempty
upper set \(U\subseteq\N^d\) is nonzero. \evidence{Kernel-verified}
\end{proposition}
\begin{proof}
In the number field \(K\) put \(S=(\sum_{j=0}^{\pi_0-1}c^ju_j)/(1-c^{\pi_0})\).
Here \(\pi_0\) is a period. Reversing the periodic indices identifies
\(S\) with \(-C_t\) in S08, equation (3), for \(q=c\).
In every complex embedding the final numerator term dominates the sum
of the absolute values of all preceding terms. Thus every conjugate
has absolute value strictly positive and at most \(1/(c-1)<1\).
Its norm has absolute value strictly between zero and one. Hence \(S\)
is not an algebraic integer. In \(\Q_p\) it is the sum of the convergent
one-dimensional series. The finite sums \(B_l=\sum_{v<l}c^vu_v\) are
algebraic integers.

Dickson's lemma gives finitely many minimal vertices \(w_1,\ldots,w_g\),
where \(g\ge1\). For nonempty \(I\subseteq\{1,\ldots,g\}\), let \(w_I\)
be the coordinatewise maximum. The intersection of the corresponding
orthants has mass \(\prod_i(S-B_{(w_I)_i})\).
Inclusion--exclusion therefore expresses \(H(U)=F(S)\), where
\[
 F(X)=\sum_{\varnothing\ne I\subseteq\{1,\ldots,g\}}(-1)^{|I|+1}
                 \prod_{i=1}^d(X-B_{(w_I)_i}).
\]
The leading coefficient is \(1-(1-1)^g=1\). A zero would make \(S\)
integral over the ring of integers of \(K\), hence over \(\Z\).
This is the \lean{PeriodicPolynomialMass}{174}{periodic mass theorem}.
\end{proof}

Fix a prime \(p\mid c\), put \(\rho=v_p(c)\), and let \(\omega\in\Z_p\) be the
root of unity of order dividing \(p-1\) congruent to \(b^{-1}\bmod p\)
(\(\omega=1\) at two). A principal unit raised to the \(p\)th power gains
at least one power of \(p\) in its difference from one, by the binomial
theorem. Thus \(b^{-c^j}\equiv\omega^{c^j}\pmod {p^{1+\rho j}}\).
The residues of \(c^j\) modulo the order of \(\omega\) are eventually
periodic, with period \(\pi_0\). The phase is \(k\bmod\pi_0\). In a fixed class of \(k\bmod\pi_0\), define
\(u_v=\omega^{c^j}\) using any late \(j\equiv k-v\bmod\pi_0\).

For example, in base two with \(c=35\), the local cycles at 5 and 7 both
have length two. At 5, for even shell index, \(u_v=(-1)^v\omega\) with
\(\omega^2=-1\). Thus \(S=\omega/36\), and in degree two the cut
\(1+1/35\le t<2\) omits only the origin. Its mass
\(S^2-\omega^2=1295/1296\) is divisible by 5. Bounded loss allows this
cancellation.

\begin{lemma}[Reflection at fixed precision]\label{refllemma}
For every \(L\ge1\), uniformly in \(n\) on late shells of this class,
\begin{equation}\label{reflection}
 c^{dk} b^{-n}R_d(n)\equiv H(U_t)
\pmod{p^L}.
\end{equation}
\evidence{Paper proof}
\end{lemma}
\begin{proof}
Reflecting each retained tuple gives its weight
\(\prod_i c^{v_i}b^{-c^{k-v_i}}\). With \(|v|=\sum_i v_i\), all terms of \(\rho|v|\ge L\) vanish at
this precision. On the remaining finite set the indices \(k-v_i\) are
positive, in the periodic range and satisfy \(1+\rho(k-v_i)\ge L\) once \(k\) is sufficiently large.
\end{proof}

\begin{lemma}[Uniform loss]\label{losslemma}
For fixed \(b,c,d,p\) with \(c>2\), the losses on all cuts \(U_t\)
are bounded independently of the shell and phase. \evidence{Paper proof}
\end{lemma}
\begin{proof}
The indicators of upper sets containing the fixed anchor form a compact
space in the product topology on \(\{0,1\}^{\N^d}\), since all defining
conditions are closed. The mass is continuous into \(\Z_p\) by finite
truncation at each precision. Its compact image omits zero, so one \(L\)
bounds the loss \(v_p(H(U))<L\) on the entire family. Take the maximum
over the finitely many phases.
\end{proof}

\begin{lemma}[Denominators of polynomial prefixes]\label{poles}
Let \(b\ge2,\ c>2\) be coprime. For integral \(Y\) of degree \(d\ge1\), \(q_n=\den R_Y(n)\) divides
\(c^{dk}\) and satisfies \(q_n\asymp n^d\) on late shells.
For rational \(P\) of degree \(d\), a fixed radix shift gives denominators
of the same order with prime support contained in a fixed set coprime to \(b\).
\evidence{Paper proof}
\end{lemma}
\begin{proof}
By Lemmas~\ref{refllemma} and \ref{losslemma}, the denominator of
the leading prefix retains \(p^{\rho dk-O(1)}\).
For an integral polynomial \(Y\) of degree \(d\), multiplication by the
leading coefficient costs a fixed order. Lower powers have denominator
dividing \(c^{(d-1)k}\). Their poles are strictly smaller for large \(k\).
Applying this at \emph{each} prime of \(c\) shows that the reduced
denominator \(q_n\) of \(R_Y(n)\) satisfies \(q_n\asymp n^d\) on late
shells and divides \(c^{dk}\).

For rational \(P\), choose \(t_0\ge0\) so \(b^{t_0}P=Y/D_0\) with integral
\(Y\) and \(D_0\) coprime to \(b\), by removing the radix primes of a common
denominator. The new prefix denominator \(Q_n\) satisfies
\(q_n\mid Q_n\mid D_0 q_n\); its prime support is contained in a fixed set coprime to \(b\).
\end{proof}

\begin{lemma}[Cancellation on blocks between insertions]\label{block}
Under the hypotheses of Theorem~\ref{main}, each nonzero integer
frequency has vanishing mean on the radix orbit of \(P(\alpha)\).
\evidence{Lean-supported paper proof}
\end{lemma}
\begin{proof}
Use the shifted polynomial \(Y/D_0=b^{t_0}P\) from Lemma~\ref{poles}.
The real error in replacing \(b^n\alpha^\ell\) by \(R_\ell(n)\) tends to zero:
every fixed tuple eventually enters, and its omitted weight is dominated
by \(c^{-\sum j_i}\), a summable majorant.

Fix \(h'\in\Z\setminus\{0\}\). Up to \(N\) there are
\(O((1+\log N)^d)\) insertion locations. Discard the first and last
incomplete portions of length at most \(B=\lfloor\sqrt N\rfloor\) each,
starting the partition at \(B\). Discard also any length-\(B\) block
containing an insertion. The discarded fraction tends to zero.
On every remaining block the rational prefix follows an exact radix orbit.
At frequency \(h'\) its initial reduced denominator is
\(Q_n/\gcd(|h'|,Q_n)\), tending uniformly to infinity and at most \(B^{3d}\)
for large \(N\). The \lean{SmoothOrbit}{207}{fixed-support short-orbit
estimate} gives a uniformly vanishing mean of \(\ee{h'b^i R_Y(n)/D_0}\)
on these blocks. This estimate states cancellation for every reduced
denominator \(q\to\infty\) supported on a fixed finite set coprime to \(b\)
when \(q\le B^C\) for a fixed \(C\). Vandehey's Corollary~1.2 gives a
stronger quantitative form. The vanishing real error transfers cancellation
to the actual orbit. Removing finitely many initial positions preserves
a limiting mean.
\end{proof}
\begin{proof}[Proof of Theorem~\ref{main}]
Lemma~\ref{block} and Weyl's criterion give uniform distribution.
A shift of \(t_0\) positions changes each prefix count by at most \(t_0\).
\end{proof}
\begin{proof}[Proof of Theorem~\ref{joint}]
For any nonzero integer vector \(h\), the polynomial \(\sum_i h_iX^{r_i}\)
has a surviving coefficient of positive degree. Theorem~\ref{main} gives
vanishing of the corresponding character mean. Multidimensional Weyl
equidistribution follows by uniform approximation of continuous functions
by trigonometric polynomials. Every half-open box has boundary of Haar
measure zero, including boxes of zero volume. For \(m=0\) the statement is trivial.
\end{proof}

\paragraph{Other conclusions of the one-axis proof.}
For a common odd prime \(p\nmid b\), products of factors
\(\alpha_{b^{m_i},p^{e_i}}\), \(m_i,e_i\ge1\), permit independent strides.
A rational combination with a unique nonzero term of largest positive
degree is normal \cite[Theorem~2]{archive} \evidence{Kernel-verified}
(\lean{StonehamPolynomial}{107}{one-axis sum}). Indeed, reflect each index
about \(\lfloor\log_p(n)/e_i\rfloor\). The weight becomes
\(p^{\sum e_i v_i}\) times a common unit modulo any fixed precision on a
late shell, since \(p^s\bmod\varphi(p^L)\) stabilizes. The inclusion--exclusion
numerator of this weighted series has value one at \(1\) and hence is
\(1\bmod(p-1)\) at \(p\). The anchor and pole comparison above apply.
For \(\alpha_{2,3}^2\) the exact late-shell denominator exponents on the four
successive ranges cut by \(3^k+3^{k-2},3^k+3^{k-1},2\cdot3^k\) are
\(2k-2,2k-3,2k-1,2k\) for \(k\ge4\)
(\lean{StonehamSquare}{251}{shell profile}) \evidence{Kernel-verified}. The initial dip at \(n=30\)
is from 81 to 9. The archived proof counts ordered pairs and treats the
missing index-zero terms separately \cite[Theorem~7]{archive}.

\section{What changes at parameter two}\label{two}
Now let \(b\ge3\) be odd and \(\alpha=\alpha_{b,2}\). The local units tend
to one, and the full one-dimensional mass is \(\sum_{v\ge0}2^v=-1\) in
\(\Q_2\). On the square cut \(3/2\le t<2\) only the origin is omitted,
so the two-dimensional mass is \((-1)^2-1=0\).

\begin{theorem}\label{square}
The number \(\alpha^2\) is normal in base \(b\).
For \(k\ge2\) and \(2^k\le n<2^{k+1}\) its shifted square prefix has
reduced denominator
\[
 \den R_2(n)=
 \begin{cases}2^{2k-2},& n<3\cdot2^{k-1},\\
               2^{k-1},& n\ge3\cdot2^{k-1}.
 \end{cases}
\]
For \(n<4\) the denominator is one. \evidence{Kernel-verified}
\end{theorem}
\begin{proof}
Set \(S_k=\sum_{j=1}^k2^{k-j}b^{-2^j}\) and \(A_k=b^{2^k}S_k\).
This odd integer satisfies \(A_k\equiv2^k-1\bmod 2^{k+2}\),
since \(b^{2^j}\equiv1\bmod2^{j+2}\) for \(j\ge1\) by repeated squaring.
On the lower half all pairs through \(k-1\) have entered. After scaling
by \(2^{2k-2}\) their contribution is \(b^{n-2^k}A_{k-1}^2\), odd.
Each remaining off-diagonal pair comes with its transpose, and its smaller
index is at most \(k-2\), giving an even total. On the upper half all
pairs through \(k\) except \((k,k)\) have entered, so
\(b^{2^{k+1}-n}R_2(n)=(A_k^2-1)/2^{2k}\).
The two factors \(A_k+1,A_k-1\) have valuations \(k,1\), respectively.
This proves the \lean{ParameterTwoSquareProfile}{155}{profile}.
The minimum denominator on successive shells tends to infinity; the block
argument of Section~\ref{mass} gives
\lean{ParameterTwoSquare}{20}{square normality}.
\end{proof}

\begin{theorem}\label{quad}
For integers \(A\ne0,B,C\) and \(D\ge1\) put
\(x=(A\alpha^2+B\alpha+C)/D\). Let \(q_k\) be the reduced denominator
of the \emph{numerator polynomial's} shifted prefix at \(n=3\cdot2^{k-1}\).
Then \(x\) is normal in base \(b\) if and only if \(q_k\to\infty\).
\evidence{Kernel-verified}
\end{theorem}
\begin{proof}
The lower-half square pole, after multiplication by \(A\), eventually
exceeds the linear pole, which is at most \(k\). The upper half has no
further insertions and multiplication by \(b\) preserves a dyadic
denominator. Growth of \(q_k\) therefore gives growth on all positions
for the numerator. Removing radix primes of \(D\) by a fixed shift gives
the same conclusion for \(x\), and the block proof applies. Conversely,
normality of \(x\) gives uniform distribution of the orbit of \(Dx\) by
multiplication on the circle. Differences of two prefix counts show that
the upper halves, whose lengths are \(2^{k-1}\), are uniformly distributed
as \(k\to\infty\). If \(q_k\) were bounded on a subsequence, all rational
prefixes on those halves would lie on a fixed finite dyadic grid.
The real errors tend uniformly to zero, so an interval between grid points
contradicts uniformity. See the \lean{ParameterTwoQuadraticOrbits}{16}{equivalence}.
\end{proof}

\begin{proposition}[A sufficient coefficient test]\label{coeff}
Let \(b\ge3\) be odd, \(\alpha=\alpha_{b,2}\), and let \(A,B,C,D,q_k\)
be as in Theorem~\ref{quad}. Put \(\beta=v_2(b^2-1)-1\).
If, with \(v_2(0)=+\infty\),
\begin{equation}\label{nonres}
 v_2(2A+B)<v_2(2A)+\beta ,
\end{equation}
then \(q_k=2^{\max(0,k-v_2(2A+B))}\) for \(k\ge2\), and
\((A\alpha^2+B\alpha+C)/D\) is normal in base \(b\).
\evidence{Lean-supported paper proof}
\end{proposition}
\begin{proof}
For \(j\ge1\) put \(\delta_j=(b^{-2^j}-1)/2^j\).
Successive squaring gives \(\delta_j\in2^\beta\Z_2\).
The unshifted linear prefix equals \(H_k-2^{-k}\), with
\(H_k=1+\sum_{j=1}^k\delta_j\). Subtracting the missing diagonal
\((2^{-k}+\delta_k)^2\) from its square shows that the upper-half
unshifted numerator polynomial, modulo \(\Z_2\), equals
\[
 -2^{-k}\bigl(2A+B+2A(\textstyle\sum_{j=1}^k\delta_j+\delta_k)\bigr).
\]
The shift is a 2-adic unit. Under \eqref{nonres} the perturbation inside
the parentheses has strictly higher valuation than \(2A+B\).
Taking the denominator gives the formula. The kernel-verified equivalence
\lean{ParameterTwoQuadraticOrbits}{16}{\texttt{normal\_iff\_upper\_denominator}}
then gives normality.
\end{proof}
Call a quadratic \emph{resonant} here when \(2A+B=0\): the coefficient
of the leading upper-half term cancels identically.
For the resonant polynomial \(X^2-2X\), put \(Z_1=2-2b^2\) and
\(Z_{k+1}=2+2b^{2^k}(Z_k-1)\). Its
\lean{Phase16StonehamBounds}{126}{exact denominator} is
\(q_k=2^{\max(0,2k-1-v_2(Z_k))}\) \evidence{Kernel-verified}.
Whether this tends to infinity for even one specified odd radix remains
open. For instance, at \(b=3\) the prefix for this polynomial at
\(n=12\) is \(7315/4-2(249399/8)=-60521\). This integer accounts for
one upper half; failure of normality requires a bounded subsequence of
upper-half denominators.

\section{Degrees certified by the dyadic mass}\label{certs}
Write \(M_d(t)=\sum_{v\in U_t}2^{|v|}\) for \(c=2\).
The \emph{Hilbert series} counts tuples by total deficit: its coefficient
of \(z^j\) is the number in the cut with \(|v|=j\).
The orthant calculation gives
\[
 \sum_{v\in U_t}z^{|v|}=\frac{K_{d,t}(z)}{(1-z)^d},
 \qquad K_{d,t}\in\Z[z],\qquad M_d(t)=(-1)^d K_{d,t}(2).
\]
These identities also define the integer mass at a real threshold outside
the basic shell. A uniform bound \(v_2(M_d(t))<L\) on \(1\le t<2\)
implies normality for every polynomial of degree \(d\):
at \(k\ge2L\) every term of weight exponent below \(L\) has positive
original indices and radix unit \(1\bmod2^L\). The leading denominator
exponent is \(dk-v_2(M_d(t))\). The lower-degree comparison and the orbit
argument of Section~\ref{mass} apply.

For later use, nonvanishing on all cuts of this shell implies such a
uniform bound. A convergent sequence of thresholds has a monotone
subsequence with limit in \([1,2]\). On approach from above the indicators
converge to the weak cut. On approach from below the strict cut equals an
actual cut a little to its left: sums of \(d\) powers \(2^{-v_i}\) cannot
approach a limit strictly from below. Indeed, from any such sequence one
can extract a subsequence on which each exponent is constant or tends to
infinity. The other summands are nonnegative, so each sum is at least
its limit. This also handles approach to 2; the endpoint 1 can only be
approached from within the shell on its right. Mass is continuous for
pointwise convergence of indicators, by truncation modulo \(2^L\).

\begin{theorem}\label{special}
Fix an odd radix \(b\ge3\) and put \(\alpha=\alpha_{b,2}\).
Every rational polynomial of degree four or of degree \(p^r\) for an odd
prime \(p\) and integer \(r\ge0\) has a normal value at \(\alpha\) in
base \(b\). \evidence{Kernel-verified}
\end{theorem}
\begin{proof}
For \(d=p^r\) use the orbit argument below with \(e=1\), where
\(M_1(y)=-1\) throughout the shell. For \(r=0\) this last identity applies
directly. In degree four the \lean{DyadicMassCells}{151}{finite grid
calculation} bounds the loss on all cuts: after removing the factor
\(2^4\) from the zero-index normalization, it is \(8,7,8\) on
\([1,3/2),[3/2,7/4),[7/4,2)\), respectively. The middle interval
joins two cells meeting at \(13/8\): their residues modulo \(2^{13}\)
are \(2^{11}\) and \(3\cdot2^{11}\), both of valuation \(11=4+7\). Reflection and the preceding
transfer give the \lean{StonehamOddDegree}{63}{prime-power} and
\lean{DyadicQuartic}{109}{quartic} theorems.
\end{proof}

\begin{proposition}\label{finite}
Fix an odd radix \(b\ge3\) and put \(\alpha=\alpha_{b,2}\).
For \(3\le d\le120\), and for \(d=128\), a finite binary partition of
\([1,2)\) certifies that \(v_2(M_d(t))\) is finite throughout.
For each pair \((e,p)\) in Table~\ref{pairs}, the residues of
\(M_e(t)\bmod p\) omit zero throughout this interval.
\evidence{Finite certificate}
\end{proposition}

\begin{table}[ht]
\centering
\begin{tabular}{cl}
\toprule \(e\)&Primes \(p\)\\ \midrule
3 & \(31,127,257,337,1801,8191,65537,131071,524287,2147483647\)\\
4 & \(8191,65537,131071,524287,2147483647\)\\
5 & \(524287\)\\ \bottomrule
\end{tabular}
\caption{Sixteen pairs with a closed digit orbit avoiding zero modulo \(p\).}
\label{pairs}
\end{table}

\begin{proof}[Verification of the certificate]
Apply the two stopping rules of Lemmas~\ref{tree} and \ref{words} below.
With maximum degree 120, the
\repo{scripts/phase16_dyadic_stoneham.py}{first verifier} checks the least
passing precision in each degree. For example \(d=3,4,15,120\) have
largest losses \(6,8,64,835\). At \(d=120\) the least loss is 716
and the first largest loss occurs at \(2157/2048\).
At \(d=128\), the \repo{scripts/phase17_dyadic_stoneham.py}{second verifier}
with its degree-128 option closes 345 leaves at precision 907.
The least and greatest losses are 776 and 906; the latter first occurs at
\(1+2^{-5}+2^{-6}+2^{-7}+2^{-17}\).

The first verifier covers seven pairs and the second the other nine;
the latter includes an orbit of 338544 vectors for \((5,524287)\).
For \((3,31)\) the reduced orbit has 57 vectors in dimension five.
Lean verifies the \lean{Phase17DiscrepancyBounds}{563}{Boolean certificate}
and \lean{Phase17DiscrepancyBounds}{579}{induction on all words}.
The identification of the reduced matrices with this mass recurrence uses
the intertwining identities checked by the exact script.
\end{proof}


\begin{theorem}\label{degrees}
Fix an odd radix \(b\ge3\) and put \(\alpha=\alpha_{b,2}\).
Every rational polynomial of degree \(3\le d\le120\), of degree 128, or
of degree \(d=ep^r\) for a pair in Table~\ref{pairs} and \(r\ge0\), is
normal at \(\alpha\) in base \(b\). This uses the finite certificate of
Proposition~\ref{finite}. \evidence{Lean-supported paper proof}
\end{theorem}
\begin{proof}
For the finite range use Proposition~\ref{finite} and
\lean{DyadicQuartic}{80}{normality from retention}. The zero-index
prefix in this implication has insertion multiplier 2, hence equals
\(2^d R_d(n)\) by shifting every coordinate by one. This changes the
denominator loss by at most \(d\).
For \(r\ge1\), rotate each of \(e\) blocks of \(p^r\) coordinates by one
place simultaneously. On each finite weight slice every nonfixed orbit
has size divisible by \(p\). Fixed tuples are constant on each block;
their free coordinates satisfy the cut at \(t/p^r\). Thus in \(\mathbb F_p[[z]]\)
the Hilbert series at \((ep^r,t)\) equals the one at \((e,t/p^r)\)
with \(z\) replaced by \(z^{p^r}\). Multiplication by
\((1-z)^{ep^r}=(1-z^{p^r})^e\bmod p\) gives a polynomial identity.
Evaluation at \(2\) and Fermat's congruence give
\(M_{ep^r}(t)\equiv M_e(t/p^r)\bmod p\), the signs agreeing as \(p\)
is odd. While \(x<1\) all deficits are positive and
\(M_e(x)=2^e M_e(2x)\). Doubling \(t/p^r\) to \(y\in[1,2)\)
therefore gives a unit multiple of \(M_e(y)\bmod p\). The certificate
excludes zero, and compactness of the cuts gives bounded dyadic loss for
each fixed degree. The \(r=0\) case is the certificate itself.
\end{proof}

\begin{corollary}\label{echelon}
Fix an odd radix \(b\ge3\) and put \(\alpha=\alpha_{b,2}\).
Let \(P_1,\ldots,P_m\in\Q[X]\) be linearly independent. Choose a basis
of their span with distinct leading degrees (the pivot degrees). Assume
the span contains no nonzero constant, equivalently that all pivot degrees
are positive. Each pivot degree must be covered by Theorem~\ref{special}
or \ref{degrees}, or equal to 2.
If 2 is a pivot degree, assume 1 is absent and that the
quadratic basis element satisfies \eqref{nonres} after clearing its
denominators. Then the vector of values \(P_i(\alpha)\) has the
half-open box limits of Theorem~\ref{joint}. Pivot degrees taken from
Theorem~\ref{degrees} use the finite certificate of Proposition~\ref{finite}.
\evidence{Lean-supported paper proof}
\end{corollary}
\begin{proof}
A nonzero integer combination of the \(P_i\) has the highest occupied pivot
degree in its expansion in this basis. If the degree is two, no lower
pivot degree exists, so the polynomial is a nonzero rational multiple of the
quadratic basis element. Inequality~\eqref{nonres} is invariant under this scaling.
Every nonzero character thus has a normal scalar projection. Apply the
Fourier argument of Theorem~\ref{joint}. The \lean{Phase17DiscrepancyBounds}{311}
{formal echelon implication} starts with integer polynomials of distinct
degrees and assumes normality for \emph{all} integer polynomials of each
such degree. The restricted quadratic pivot and the change of rational
basis above are paper steps.
\end{proof}
For example \((\alpha^2+\alpha,\alpha^7)\) satisfies the criterion at
\(b=3\). The span of \(\alpha,\alpha^2\) at the polynomial level includes
\(X^2-2X\); the separate proofs of normality of its two coordinates do
not decide their joint distribution.

\subsection*{How to read the certificates}
At a binary prefix of \(t\), store an integer constant and a list of
integer coefficients indexed by pairs \((\ell,n)\). A pair denotes the
unknown function \(M_\ell(n+f)\) of the unread fraction \(0\le f<1\).
Start with the single coefficient one at \((d,1)\) for \(d\ge2\).
The following update reads one bit and combines equal successor states.
For a state \(M_\ell(n+f)\), \(0\le n<\ell,\ 0\le f<1\), write
\(2f=\epsilon+f'\) with first binary digit \(\epsilon\in\{0,1\}\). Choosing the
\(m\) coordinates of deficit zero and subtracting one from all other
deficits gives
\[
 M_\ell(n+f)=\sum_{m=0}^n \binom{\ell}{m}2^{\ell-m}
                  M_{\ell-m}(2(n-m)+\epsilon+f').
\]
If the successor argument is at least its dimension \(\ell'\), the entire
orthant has entered and its mass is \((-1)^{\ell'}\). Add this term to the
constant (it is \emph{absorbed}). A transition to any other state is
\emph{live} and retains dependence on the unread digits.

For the zero continuation, the same recurrence computes the integer
masses \(M_\ell(n)\) exactly: \(M_\ell(0)=0\) for \(\ell>0\), and \(n\ge \ell\)
gives \((-1)^\ell\). For \(0<n<\ell\) the \(m=0\) term doubles \(n\);
every \(m>0\) term reduces the dimension. Thus the recursion terminates.

\begin{lemma}[Binary interval certificate]\label{tree}
At a prefix node let \(C\) be the absorbed constant and \(w_{\ell,n}\)
the live coefficients. If \(v_2(C)<v_2(w_{\ell,n})+\ell-n\) for every live
state, the loss on the entire half-open interval is \(v_2(C)\).
With no live states the mass equals \(C\). \evidence{Paper proof}
\end{lemma}
\begin{proof}
The update partitions the ordered tuples by their zero deficits; every
remaining deficit contributes a factor two upon subtracting one.
Thus it preserves the represented function exactly. Since \(m\le n\)
every term from \((\ell,n)\) has factor \(2^{\ell-n}\); this gives the test
by addition in \(\Z_2\). More explicitly, at precision \(L\) put
\(g=\min(L,\min_{\ell,n}(v_2(w_{\ell,n})+\ell-n))\), using \(+\infty\)
for an empty inner minimum. All continuations agree with the value at
the left endpoint modulo \(2^g\). Stop if this residue is nonzero, or
if \(g=L\), in which case a zero residue fails at this precision.
There is a second stopping rule. A sum of \(d\)
binary powers has at most \(d\) ones after carrying. If the read prefix
already contains \(d\) ones including the integer digit, any threshold
strictly above its zero continuation and inside its interval would have
an additional one. Hence the value at that continuation serves the whole
interval. Use trailing zeros at dyadic endpoints. Starting at the root,
stop by either rule or split into both sons. Every split gives disjoint
half-open intervals covering the parent. Thus a finite tree closed with
nonzero leaf values certifies every real threshold, including endpoints.
At a fixed precision coefficients divisible by that power of two can be
discarded; each live update gains at least one bit.
\end{proof}

\begin{lemma}[Odd-prime word certificate]\label{words}
Reduce the same two updates modulo an odd prime \(p\). If every reachable
coefficient row gives a nonzero value at the zero continuation, then
\(M_e(t)\ne0\bmod p\) throughout the shell. \evidence{Paper proof}
\end{lemma}
\begin{proof}
Append the absorbed constant as a coordinate; evaluate a row against the
column of integer-argument masses, with a one for that coordinate.
Starting from \((e,1)\), adjoin both successor rows until no new row
appears. The dimension is finite (\(1\le \ell\le e,\ 0\le n<\ell\)), so
closure terminates over \(\mathbb F_p\). Induction on word length puts
every finite word in this orbit. Every real threshold has either a
terminating expansion or a prefix containing its \(e\)th one. In the
second case the zero continuation has the same mass by Lemma~\ref{tree}.
\end{proof}

\section{The least orbit order in every degree}\label{leastsec}
A permutation orbit of deficits is described by multiplicities
\(m_i\), only finitely many positive, with \(\sum m_i=d\).
An orbit consists of all ordered tuples obtained by permuting the deficits.
Write
\[
 K(m)=\sum_i m_i2^{-i},\qquad
 \Phi(m)=\sum_i\bigl(i m_i-v_2(m_i!)\bigr),\qquad
 F(d)=\min_{K(m)<2,\,\sum m_i=d}\Phi(m).
\]
The orbit contributes \(2^{\sum i m_i} d!/\prod m_i!\) to \(M_d(t)\)
when \(K(m)\le t\). Its orbit order, meaning the 2-adic valuation of this integer, is \(v_2(d!)+\Phi(m)\).

\begin{theorem}\label{least}
Fix an odd radix \(b\ge3\) and put \(\alpha=\alpha_{b,2}\).
For all \(d\ge1\),
\[
 F(d)=\sum_{k=1}^d\lfloor\log_2 k\rfloor .
\]
Let \(2^r-1\le d\le2^{r+1}-2\) and \(j=d-2^r+1\).
For \(d\ge2\) the minimum valuation of \(M_d(t)\) on \(1\le t<2\) is
\(v_2(d!)+F(d)\). It occurs precisely on
\[
 J_d=
 \begin{cases}
 [2-2^{1-r},\,2-2^{-r}),& j>0,\\
 [2-2^{1-r},\,2),&j=0.
 \end{cases}
\]
In particular \(M_d(t)\ne0\) there. \evidence{Paper proof}
\end{theorem}
\begin{proof}
Denote the sum of logarithmic floors by \(G\). For each \(D\ge0\) put
\[
 N_D=\sum_{i\le D}m_i,\qquad A_D=\sum_{i\le D}2^{D-i}m_i,\qquad
 S_D=\sum_{i>D}\lfloor m_i/2^{i-D}\rfloor .
\]
Counting the levels below each \(i\) gives
\(\sum i m_i=\sum_D(d-N_D)\), and counting the powers of two up to
each \(k\) gives \(G=\sum_D\max(0,d+1-2^{D+1})\).
Since \(m_i<2^{i+1}\), Legendre's factorial formula gives
\(\sum_i v_2(m_i!)=\sum_D S_D\), including \(i=0\). Subtraction yields
\[
 \Phi(m)-G=\sum_{D\ge0}
 \bigl[\min(d-N_D,2^{D+1}-1-N_D)-S_D\bigr].
\]
Each summand is nonnegative. Indeed \(S_D\le d-N_D\), and
\(A_D+S_D\le 2^D K(m)<2^{D+1}\), with \(A_D\ge N_D\).

Suppose equality holds and the largest occupied index is \(I\ge1\).
For \(D<I\), \(S_D<d-N_D\), so the second bound must be tight, forcing
\(A_D=N_D\). At \(D=I-1\) this leaves only levels \(I-1,I\) with
multiplicities \(x,y\) satisfying \(y>0,\ x+\lfloor y/2\rfloor=2^I-1\).
Conversely, any such pair has \(K(m)<2\) and attains equality: for
\(D=I-1-m\), \(m\ge1\), the binary digits of \(x,\lfloor y/2\rfloor\)
are complementary through \(I\) bits. The sum of their quotients rounded
down after division by \(2^m\) is \(2^{I-m}-1\), exactly the second
bound. For \(D\ge I\) the summands are zero.

For \(y\) even one obtains \(I=r,\ (x,y)=(2^r-1-j,2j)\) when \(j>0\).
For \(y\) odd and \(j>0\) one obtains
\(I=r,\ (x,y)=(2^r-j,2j-1)\). These exhaust the choices since
\(d=2^I-1+\lceil y/2\rceil\) and \(0\le\lfloor y/2\rfloor\le2^I-1\).
At \(d=2^r-1\) the sole choice has all entries at \(r-1\)
(including \(d=1,I=0\)). This proves the minimum formula.
For \(j>0\) the two minimizing orbits enter at \(2-2^{1-r}\) and
\(2-2^{-r}\), in that order. Exactly one odd multiple of
\(2^{v_2(d!)+F(d)}\) is present on \(J_d\). Before it there is none;
after the second entry there are two, whose sum is even at that scale.
For \(j=0\) the sole minimizer enters at the left endpoint of \(J_d\).
All other contributions have greater order, and the series converges
2-adically. This proves the valuation assertion.
\end{proof}
For \(d=2\), the interval is \([1,3/2)\). The vanishing on the other half
shows why a least-order calculation on one interval cannot by itself
prove a whole-shell denominator bound. Whether \(M_d\) is nonzero on
the entire shell for every \(d\ge3\) remains open.

\paragraph{Open problems and the real shell parameter.}
Whole-shell nonvanishing is stronger than the following sufficient test.
\begin{proposition}[A test on open intervals]\label{plateaus}
For \(d\ge1\), let \(D_*=\{0\}\cup\{2^{-j}:j\ge0\}\) and let \(dD_*\)
be the set of sums of \(d\) elements of \(D_*\). The complementary open
intervals in \((1,2)\) carry full Lebesgue measure on the shell, and
\(M_d\) is constant on each. Thus its distribution under uniform sampling
of \(t\) is purely atomic, and its zero set has measure zero exactly when
every such interval has nonzero mass. Under this condition every degree-\(d\)
rational polynomial in \(\alpha_{b,2}\) is normal in each odd base \(b\ge3\).
\evidence{Paper proof}
\end{proposition}
\begin{proof}
The finite Cartesian power of the compact countable set \(D_*\) is
compact and countable. Its image under addition has the same properties
and contains all tuple thresholds. Hence the indicators of the cut, and
therefore the mass, are constant on each complementary interval.
There are countably many intervals; each has positive length, and the
exceptional points have measure zero. This proves the distribution and
zero-set assertions.

Suppose every open plateau has nonzero mass. For fixed \(L\), the failed
set \(E_L=\{t\in[1,2):M_d(t)\equiv0\bmod 2^L\}\) is a finite union of
intervals since only \(|v|<L\) matters. By continuity of measure from
above, \(|E_L|\to0\). At \(n=2^kt\) the finitely many retained deficits
see \(b^{-2^{k-v_i}}\equiv1\bmod 2^L\) uniformly for large \(k\).
Outside \(E_L\) the leading prefix therefore retains denominator exponent
at least \(dk-L+1\). After clearing coefficients, the leading coefficient
costs a fixed valuation and all lower powers have depth at most \((d-1)k\).
Radix-prime factors of a rational coefficient denominator are removed by
the shift in Lemma~\ref{poles}.

For any fixed denominator cap, all sufficiently late positions below
the cap must therefore lie in failed cells. Through a cutoff \(N\) in
shell \(K\), their number is at most
\(\sum_{k\le K}2^k|E_L|+O_L(K+1)+O_L(1)\), by counting integers in
each interval, even on the partial final shell. As \(N\ge2^K\), their
upper density is at most \(2|E_L|\), hence zero. In the block proof of
Lemma~\ref{block} every insertion-free block has constant reduced
denominator. For a fixed frequency the additional reduction is bounded
by its absolute value. The blocks below any fixed cap thus have total
length \(o(N)\); the uniform short-orbit estimate on the remaining
blocks tends to zero as the cap increases. The insertion count, upper
denominator bound and vanishing real tail are as in that lemma.
\end{proof}
The values on the open plateaus in the unresolved degrees still have to
be determined. In degree two the mass actually vanishes on \([3/2,2)\).
For comparison, the synthetic cubic schedule
\(x=\sum_{i\ge1}3^{-v_3(i)}2^{-2i^3}\) admits a curvature argument for
normality of \(x^2\): isolated pairs on a curved grid produce jumps
with unit coefficient \(2\bmod 3\). Dyadic Stoneham thresholds are sums
of powers of two and have the plateau distribution just proved. That
curvature argument does not apply to \(\alpha_{b,2}^2-2\alpha_{b,2}\).
Its normality for any specified odd radix remains open, as does
whole-shell nonvanishing of \(M_d\) for all \(d\ge3\).

\section{Discrepancy through an arbitrary digit}\label{ratesec}
For the ordinary case \(c>2\), Mira Cult's S08, Theorem 1.1(b)
\cite{mira2}, already gives an effective power saving for the vector of
powers; Corollaries 1.2--1.3 identify the ordinary and polynomial-vector
cases. The scalar refinement here is the explicit exponent \(\delta_d\).
The cases at parameter two require separate denominator estimates.

Let \(D_N(x;b)=\sup_I|N^{-1}\sum_{n<N}\mathbf1_I(\{b^nx\})-|I||\),
where the supremum is over half-open subintervals of \([0,1)\).
Write \(\Lambda_N=1+\log N\).
\begin{theorem}\label{rate}
If \(c>2\) under the hypotheses of Theorem~\ref{main}, then for \(d=\deg P\ge1\)
and every \(N\ge2\),
\[
 D_N(P(\alpha_{b,c});b)\le C(b,c,P)N^{-\delta_d},\qquad
 \delta_d=\frac{1}{2(2^{d+1}-2)}.
\]
The constant is computable. For \(c=2\) and odd \(b\ge3\) the same estimate holds in
degree one, in odd prime-power degrees, and for quadratics satisfying
\eqref{nonres}. In the indicated degrees \(1,2,3\) the respective
bounds are \(O(N^{-1/2}\Lambda_N^2)\), \(O(N^{-1/6}\Lambda_N^{3/2})\) and
\(O(N^{-1/14}\Lambda_N^{5/4})\). \evidence{Lean-supported paper proof}
\end{theorem}

\begin{proof}
We give the estimates for the last incomplete shell as well as whole
shells. Use the fixed shift and \(D_0\) of Section~\ref{mass}. If a tuple
is omitted at \(n\), one coordinate has \(c^{j_i}\ge(n+1)/d'\) for its
degree \(d'\ge1\). Summing that geometric tail and the unrestricted weights
in other coordinates gives a total real error \(O(1/(n+1))\).
On \(X\le n<2X,\ X=2^u\), the rational denominators obey
\(X^\lambda\ll Q_n\ll X^d\), with \(\lambda=d\) by the pole bounds above,
except that \(\lambda=1\) for dyadic quadratics by Section~\ref{two}.
All implicit constants below may depend on the fixed input.

For \(1\le T\le X\), split \([X,X+T)\) at its distinct interior
insertions, retaining each insertion at the start of its piece.
For each \(\gamma>0\) the lengths satisfy
\begin{equation}\label{length}
 \sum_i \ell_i^\gamma\ll_{d,\gamma} X^\gamma.
\end{equation}
Indeed, starts of pieces of length in \([U,2U)\), \(U=2^j\), are
\(U\)-separated. Except for \(X\) they are sums of at most \(d\) powers
of \(c\). Discard summands below \(U\); the new sum is within \(dU\)
of the start and there are \(O_d((1+u-j)^d)\) possible truncated
tuples. Each accounts for at most \(d+1\) separated starts.
Weighting by \((2U)^\gamma\) and summing in \(u-j\) proves \eqref{length}.
For \(d=1,\gamma=0\) there is at most one interior insertion. For
\(\nu\ge1\) one has \(\sum_i \ell_i^\nu\le X^\nu\).

Apply Lemma~\ref{vandestep} at the reduced modulus of each piece and
frequency. The sum may start at zero by replacing its numerator \(u\)
by \(ub^{-1}\bmod q\).
For frequency \(h\ge1\), \(q=Q/\gcd(h,Q)\) on a piece of starting
denominator \(Q\). The negative exponent costs at most a bounded factor
after harmonic averaging, apart from the usual logarithm:
\[
 \sum_{h\le H}\frac{\gcd(h,Q)^{a_k}}h
 \le (1+\log H)\sum_{l\mid Q}l^{a_k-1}
 \le (1+\log H)\prod_{p\mid cD_0}(1-p^{a_k-1})^{-1}.
\]
The first inequality enlarges each gcd class to all multiples of \(l\).
In the Erd\H os--Tur\'an inequality
\cite[Thm.~2.2(9), pp.~529--530]{bc}, multiplied by \(T\), the real tail costs \(O(H)\),
since the phase error costs \(O(hT/X)\) at frequency \(h\).
With \(H=\lfloor\sqrt X\rfloor,\ k=d-1\), \eqref{length} gives
uniformly in the terminal length
\[
 T D_{X,T}\ll \sqrt X+
 (1+\log X)^{1+2^{-k}}\bigl(X^{d a_k+\gamma_k}+X^{\nu_k-\lambda a_k}\bigr).
\]
Here \(D_{X,T}\) denotes discrepancy on that shell prefix.
The \lean{QuantitativeStonehamSupport}{49}{finite transfer} takes both
analytic inequalities as hypotheses.

To check the margin in the power saving use the recurrence in
Lemma~\ref{vandestep}. Set \(s_i=(1-\gamma_i)/a_i,\ t_i=(\nu_i-1)/a_i\). Substitution gives
\(s_0=2,\ t_i=s_i-4+2^{1-i}\) and
\(s_{i+1}=s_i+1-a_i t_i\).
Since \(\nu_i\ge1\), \(s_i\le i+2\); \(t_0=t_1=0\) and \(t_i\le i-1\)
for \(i\ge2\). Each finite partial sum of \(a_it_i\) is strictly below
\(\sum_{i\ge2}(i-1)2^{-i-1}=1/2\).
Thus \(s_k>k+3/2\) and \(d-t_k\ge1\).
For \(\lambda=d\), both powers in the last estimate are at most
\(1-\eta\) with \(\eta=\min(1/2,a_k(s_k-d),a_k(d-t_k))>a_k/2\).
For \((d,\lambda)=(2,1)\), the exponents are both \(5/6\).
Substitution in degrees one to three yields the stated log-refined bounds.
Summing the \emph{count} errors over complete dyadic shells and the last
partial shell costs a convergent geometric factor, since \(1-\eta\ge1/2\).
Divide by \(N\), absorb the fixed shift error and use \(\eta>\delta_d\)
to absorb the logarithm.

For effectivity, test successive precisions \(L\) in \eqref{reflection}.
At each prime, cycle phase and precision, only the finitely many
\(|v|\le(L-1)/v_p(c)\) can contribute. Test \(t=1\) and every jump in
\((1,c)\), including all tuples tied at a jump. The compactness bound
ensures termination. The preperiod, unit residues and threshold
\(k-v_i\ge\max(1,\text{preperiod},\lceil(L-1)/v_p(c)\rceil)\)
are computable. The same search terminates on the dyadic cuts of the
odd prime-power degrees. For a nonresonant quadratic use the explicit
valuation test. The remaining products, geometric bounds and constants
of the two analytic estimates are effective.
\end{proof}

\begin{lemma}[Vandehey's estimate on reduced supports]\label{vandestep}
Fix \(b\ge2\) and a finite set \(\mathcal P\) of primes coprime to \(b\).
For each \(k\ge0\) there are effective constants \(A_k,B_k\) depending
only on \(b,\mathcal P,k\) such that for every \(q\ge1\) supported on
\(\mathcal P\), \((u,q)=1\), and integer \(T\ge1\),
\[
 |S(u,q,T)|:=\left|\sum_{n=1}^T \ee{ub^n/q}\right|
 \le (A_k q^{a_k}T^{\gamma_k}+B_k q^{-a_k}T^{\nu_k})
                 (1+\log q)^{2^{-k}} .
\]
Here \(a_k=(2^{k+2}-2)^{-1}\), \((\gamma_0,\nu_0)=(0,1)\).
Writing \(a,\gamma,\nu\) for the current values, the next exponents are
\[
 a'=\frac a{2(1+a)},\quad
 \gamma'=\frac{1+\gamma+a\nu}{2(1+a)},\quad
 \nu'=\frac{1+\nu}{2}+\frac{a(1+\gamma-\nu)}{2(1+a)} .
\]
\evidence{Paper proof}
\end{lemma}
\begin{proof}
The estimate and exponents are those of \cite[preprint, Thm.~5.1,
pp.~13--14]{vandehey}. In equation (12) of the preprint the multiplier
is the product of \emph{all} ambient primes. If a prime is absent from
the current modulus, the resulting \(q'\) need not divide \(q\).
Equation (11) itself introduces smaller supports. We prove the
induction using the radical of the current modulus throughout.

The depth-zero estimate is the complete-period and incomplete-period
bound of \cite[preprint, Lemma~2.3, eq.~(7)]{vandehey}: its terms are
\(\sqrt q\) and \(O_{b,\mathcal P}(T/\sqrt q)\), both multiplied by \(1+\log q\), with effective
constant. The case \(q=1\) at every depth follows from the trivial bound
\(T\). The recursions preserve \(0\le\gamma<1\), \(1\le\nu<2\) and
\(\nu-\gamma\le1\): the new sum is \(1+(\gamma+\nu)/2<2\),
the new difference is \(a/(1+a)+(1-a)(\nu-\gamma)/(2(1+a))\le1\),
and the formula gives \(\nu'\ge1\).

Assume the bound at depth \(k\) for \emph{all} supports in \(\mathcal P\).
If \(T^{1+\gamma-\nu}\ge q\), the new second term is at least
\(T^{(1+\nu)/2}\ge T\). Otherwise put
\[
 W=q^{a/(1+a)}T^{(1+\gamma-\nu)/(1+a)}=q^x,\qquad 0<x<1.
\]
For each \(p^e\parallel q\) take \(v_p(q')=1+\lfloor xe\rfloor\),
raised to at least two if \(p=2\) and \(e\ge2\). All other valuations
are zero. Each exponent is at most \(e\). Multiplying the prime-power
bounds gives the replacement for (12)--(13):
\[
 q'\mid q,\qquad \mathrm{rad}(q')=\mathrm{rad}(q),\qquad
 W\le q'\le 2\,\mathrm{rad}(q) W.
\]

Let \(\tau\) be the multiplicative order of \(b\bmod q'\), put
\(g=\gcd(q,b^\tau-1)\) and \(D=q/g\).
Then \(q'\mid g\) and \(g\ll_{b,\mathcal P}\tau\le q'\).
Indeed \(v_p(b^t-1)\le C_p+v_p(t)\) whenever \(p\mid b^t-1\).
For odd \(p\) apply the binomial theorem to powers of \(b^{\mathrm{ord}_p(b)}\):
a \(p\)th power raises the valuation of the difference from one by one,
and a power prime to \(p\) preserves it. At two use \(b^2\equiv1\bmod8\)
for even \(t\) and \(v_2(b^t-1)=v_2(b-1)\) for odd \(t\).
Multiplication over the finite support gives an effective constant.
Moreover \(\gcd(q,b^{i\tau}-1)=g\gcd(i,D)\): at any unsaturated prime
of \(g\) the same binomial calculation adds \(v_p(i)\); at two, if
unsaturated, \(4\mid q'\) guarantees \(b^\tau\equiv1\bmod4\).

Split the original sum by \(n\bmod\tau\), apply Cauchy--Schwarz and
expand the squares. A difference \(i\tau\) gives a sum of length
\(T-i\tau\) with numerator \(u(b^{i\tau}-1)\) over \(q\).
Thus the diagonal costs \(\tau T\), and the off-diagonal is bounded by
\(2\tau\) times the sum of the absolute values for \(1\le i<T/\tau\).
Their reduced moduli are \(D/\gcd(i,D)\). For \(Y>0\),
\[
 \sum_{1\le i<Y}\gcd(i,D)^a
 \le Y\sum_{l\mid D}l^{a-1}
 \le Y\prod_{p\in\mathcal P}(1-p^{a-1})^{-1}.
\]
The first inequality counts multiples of each divisor. Applying the
induction hypothesis at these moduli, using \(D\asymp q/W\), gives
\[
 |S|^2\ll WT+
 \big[(q/W)^a T^{1+\gamma}+(q/W)^{-a}T^{1+\nu}\big]
 (1+\log q)^{2^{-k}}.
\]
By definition of \(W\), the first bracketed term equals \(WT^\nu\)
and absorbs \(WT\). Substituting \(W\) and taking square roots gives
\(q^{a'}T^{\gamma'}\) and \(q^{-a'}T^{\nu'}\) with the asserted
logarithm. Every constant above is computable from the old constants,
the local valuations and finite Euler products. Taking their maxima
over the finitely many subsets of \(\mathcal P\) at each depth makes
the induction uniform even when a prime disappears.
\end{proof}

\section{Formal verification and data availability}
The \repo{notes/phase18/verification.md}{checkpoint audit} records Lean
4.36.0 kernel replay with propositional extensionality, classical choice
and quotient soundness. The independent
\lean{StonehamStatements}{31}{public statement} uses the series starting
at one. The proof's zero-index axis at dilation \(c\) equals
\(c\alpha_{b,c}\), term by term after reindexing.
The accompanying verification map distinguishes complete declarations,
implications and finite certificates. In particular the all-degree
least-order formula is a written proof. Exact Python commands, parameters
and observed outputs are given alongside the source. The superseded
\repo{papers/01-stoneham-polynomials}{first edition} retains the one-axis
product proofs, the binary square's initial denominator table and the
worked figures.

The following names identify the complete kernel declarations for the
statements marked above. The two profile names include the namespace to
distinguish the squares.
\begin{center}\small
\begin{tabular}{@{}ll@{}}
\toprule Statement & Lean declaration \\ \midrule
A, B & \texttt{coprime\_polynomial\_normal}, \texttt{coprime\_joint\_powers}\\
Radix 48 & \texttt{square\_nonnormal\_48}\\
Proposition~\ref{massprop} & \texttt{periodic\_upper\_mass\_nonzero}\\
Common prime & \texttt{same\_prime\_sum\_normal}\\
Binary square & \texttt{StonehamSquare.exact\_profile}\\
Theorem~\ref{square} & \texttt{ParameterTwoSquare.exact\_profile}, \texttt{square\_normal}\\
Theorem~\ref{quad} & \texttt{normal\_iff\_upper\_denominator}\\
Resonant prefix & \texttt{quadratic\_upper\_den\_eq\_quadZ}\\
Theorem~\ref{special} & \texttt{odd\_prime\_power\_degree\_polynomial\_normal}\\
 & \texttt{quartic\_polynomial\_normal}\\ \bottomrule
\end{tabular}
\end{center}
The \href{https://github.com/eldernos/MathematicsResearch/blob/c2c9da683666b3b9368e5d4b5cb858caccf528f5/notes/phase19/nonlinear.md}{nonlinear schedule note}
gives the written curvature proof for the synthetic cubic schedule and
the real-shell comparison in Proposition~\ref{plateaus}.

\section{Use of AI tools}
This manuscript was produced in an AI-assisted research project. Large language models proposed conjectures and proof strategies, wrote the proofs and the text, implemented the exact computations and Lean 4 formalizations, searched the literature and reviewed drafts across several model providers. The repository owner directed the project and is responsible for its release. The Lean kernel checked the declarations listed in the verification map. The written proofs, literature comparisons and priority statements have been checked only by AI models from different providers. They have not been checked by an independent human expert. Cross-model review is not peer review.
\bibliographystyle{plainurl}
\bibliography{references}
\end{document}
